% ============================================================
% 数学\线段图\五年级.tex
% 本文件由 main.tex 通过 \input 引入，请勿单独编译
% ============================================================


\section{解答：行程问题相遇点按比例标位}

\vspace{0.4cm}

\noindent\textbf{5. 下面是从甲地到乙地的公路示意图，正中间是服务区。}

\vspace{0.4cm}
\begin{center}
\begin{tikzpicture}[>=Stealth, line width=1pt] % 整体思路：先画甲地、服务区、乙地三点共线的示意图，为后面的“相遇点位置”提供空白底图
  % 线段 (总长 14，代表 168 千米的题目空图)
  \draw (0,0) -- (14,0); % 画一条长度为 14 的水平线段，表示从甲地到乙地的公路
  % 甲地
  \fill (0,0) circle (3pt); % 在左端点画实心圆点表示甲地
  \node[below=8pt] at (0,0) {甲地}; % 在甲地下方标注名称
  % 服务区
  \fill (7,0) circle (3pt); % 在线段中点画服务区位置
  \node[below=8pt] at (7,0) {服务区}; % 在服务区下方标注
  % 乙地
  \fill (14,0) circle (3pt); % 在右端点画乙地实心点
  \node[below=8pt] at (14,0) {乙地}; % 在乙地下方标注名称
\end{tikzpicture}
\end{center}
\vspace{0.4cm}

\noindent 一辆小客车从甲地出发开往乙地，每小时行 120 千米；一辆大客车从乙地出发开往甲地，每小时行 90 千米。两车同时出发，经过 0.8 小时相遇。\\
\textbf{（1）在图上标出两车相遇地点的大概位置。}

\vspace{0.8cm}

\begin{center}
\begin{tikzpicture}[>=Stealth, line width=1pt] % 整体思路：沿用上一幅公路线段，在相同基准上再加出蓝色“相遇点”，让读者看到相遇位置略偏向乙地一侧

  % 线段作图基底 (总长度 14，代表 168 千米的解答落点图)
  \draw (0,0) -- (14,0); % 画公路基准线

  % 原始标记点
  % 甲地
  \fill (0,0) circle (3pt); % 甲地点
  \node[below=8pt] at (0,0) {甲地}; % 甲地标签

  % 服务区在中点
  \fill (7,0) circle (3pt); % 服务区点，仍在中点
  \node[below=8pt] at (7,0) {服务区}; % 服务区标签

  % 乙地极端点
  \fill (14,0) circle (3pt); % 乙地点
  \node[below=8pt] at (14,0) {乙地}; % 乙地标签

  % 解答新加的相遇点 (严格按 96:72=4:3 比例落点于坐标刻度 8)
  \fill[blue] (8,0) circle (3.5pt); % 在 x=8 处加蓝色圆点，表示相遇位置
  \node[below=8pt, blue, font=\bfseries] at (8,1) {相遇点}; % 在相遇点上方偏一点放蓝色加粗标签

\end{tikzpicture}
\end{center}

\vspace{0.6cm}
{\small\color{gray}
\textbf{解析：}\\
1. \textbf{计算各自行驶的路程}：\\
\indent 小客车（从甲地出发向右开）行驶路程：$120 \times 0.8 = 96$（千米）。\\
\indent 大客车（从乙地出发向左开）行驶路程：$90 \times 0.8 = 72$（千米）。\\
2. \textbf{计算总路程及中点真实位置}：\\
\indent 甲乙两地总路程为相遇之和：$96 + 72 = 168$（千米）。\\
\indent 服务区在公路的正中间，说明服务区距离原点甲地绝对距离为：$168 \div 2 = 84$（千米）。\\
3. \textbf{确定相遇点并按严谨数学比例作图}：\\
\indent 在相遇时，小客车已经行驶了 $96$ 千米。因为 $96 > 84$，说明相遇点肯定越过了正中间的“服务区”，继续向乙地方向偏离了 $12$ 千米。\\
\indent 在作图代码中，依据最新指示**“确保线段真实比例缩放”**的要求，我将整条 $168$ 千米的行程等比放大投影到了长度为 $14$ 的作图坐标系横轴上（即缩放比例为 $12:1$）。甲地位处坐标极限原点 $0$，服务区在黄金中点 $7$ 处。而根据比例映射关系 $168 : 96 = 14 : x$，相遇点位置的绝对理论坐标精确计算出 $x = 8$。\\
\indent 因此，图中带着蓝色粗体提示的**相遇点**，已被我严丝合缝地点在了横坐标刻度正好为 $8$ 的位置，这样的出图落点不存在哪怕一毫米的主观估算误差，是拥有绝对精确参考价值的标位。
}


\vspace{0.6cm}

{\small\color{blue!70!black}
\textbf{绘图基本原理与步骤：}\\
\textbf{1. 坐标定位与几何构建}：根据题意中的已知长度和角度，使用 \texttt{coordinate} 定义关键顶点坐标，通过 \texttt{draw} 指令按序连接成完整几何图形。线宽、颜色参数区分原图（黑色）与答案（红色 \texttt{red!70!black}）图层。\\
\textbf{2. 辅助量标注}：利用 \texttt{node} 的方向锚点（\texttt{above}、\texttt{below}、\texttt{left}、\texttt{right}）将角度、边长、面积等数据标注在图形周围，避免与线条或其他标注重叠。若涉及角弧则使用 \texttt{arc} 或 \texttt{pic[angle]} 命令渲染。\\
\textbf{3. 虚实线分层}：题干已知条件用实线绘制，解答新增的辅助线或操作痕迹用 \texttt{dashed} 虚线或彩色加粗线标识，确保读者一眼区分原有信息与作答内容。
}


%% ==============================
%% 第2题：体育用品商店购买问题——单价、数量的关系
%% ==============================


